\small
\noindent\textbf{Abstract}
%Abstract below

This study presents a comparative evaluation of the Ad hoc On-Demand Distance Vector (AODV) and Optimized Link State Routing (OLSR) protocols in a Mobile Ad Hoc Network (MANET) using the NS-3 simulator. The purpose of the research is to examine how reactive and proactive routing strategies perform under the same simulation conditions and to determine their relative strengths across key network performance metrics. A repeatable simulation scenario was designed in NS-3, and FlowMonitor was used to extract quantitative results from multiple simulation runs. The analysis focuses on packet delivery ratio, throughput, end-to-end delay, jitter, and packet loss, as these metrics provide a direct basis for assessing routing effectiveness in dynamic wireless environments.

The study is motivated by the fact that MANET routing performance is highly dependent on scenario configuration, including mobility, traffic behaviour, and network scale, meaning that controlled comparative testing remains necessary. In addition to metric extraction and graphical visualisation, AI-assisted analysis was incorporated as a secondary interpretive layer to help identify trends, organise comparative observations, and strengthen the clarity of result discussion. The AI component was applied only after simulation outputs had been generated and did not modify the original measured data.

The contribution of the study lies in providing a structured and reproducible workflow for comparing AODV and OLSR using NS-3 and FlowMonitor while also demonstrating a careful and transparent use of AI in simulation-based network analysis. The research therefore combines protocol comparison, repeatable experimentation, and assisted interpretation in a form intended to reflect both technical implementation and academic research rigour.
